\paragraph{消費者物価水準目標（CPLT）}
\label{sec:chapter5_a_shock_policies_cplt}
CPLTの式については
\ref{sec:chapter5_beta_shock_policies_cplt}
を参照のこと。
\begin{enumerate}
    \item
    \label{itm:chapter5_a_shock_policies_cplt_c_H_W}
        \(c_t^{H\to W}\)：
        総消費指数\(c_t^{H\to W}\)の図\ref{fig:chapter5_a_shock_graphs_c_H_W}を見ると、CPLTの軌道は消費者物価インフレ目標（IT-CPI）や生産者物価水準目標（PPLT）と同様に、ショック直後の初期において極端な大暴落を引き起こしていることがわかる。
        総消費水準目標（CLT）や名目総消費水準目標（NCLT）が落ち込みをほぼゼロに抑え込んでいるのとは対照的であり、全政策の中で最も大きくマイナス方向へ乖離している。
        消費の減少は厚生にマイナスに働くため、この前半における深刻な消費の低下がCPLTの厚生を大きく悪化させる最大の要因となっている。
        また需要ショック（\(\beta\)ショック）時に見られたような後半のプラス方向への巨大な乖離（オーバーシュート）も存在せず、深く落ち込んだ状態から極めて長い時間をかけて定常状態へ回帰するだけの軌道となっているため、消費増加による厚生の改善効果も得られていない。
    \item
    \label{itm:chapter5_a_shock_policies_cplt_y_H}
        \(y_t^H\)：
        生産\(y_t^H\)の図\ref{fig:chapter5_a_shock_graphs_y_H}を見ると、CPLTの軌道は総消費指数\(c_t^{H\to W}\)と同様に、消費者物価インフレ目標（IT-CPI）や生産者物価水準目標（PPLT）と極めて類似した極端な大暴落を見せていることがわかる。
        ショック直後の前半において、CPLTは全政策の中で最も大きくマイナス方向へ落ち込んでいる。
        生産の低下は労働の減少を意味するため、この前半部分においては労働の非効用を減らして厚生にプラスに働く効果が他の政策よりも大きくなっている。
        しかしながら、前述した深刻な消費の大暴落によるマイナス効果や、後述する価格分散の増大によるマイナス効果を補うほどの改善には至っておらず、結果として全体的な厚生の評価を低迷させる結果となっている。
        また需要ショック（\(\beta\)ショック）時に見られたような中盤以降の巨大なオーバーシュートも存在せず、深い谷底から極めて長い時間をかけて定常状態へ回帰するだけの軌道を描いている。
    \item
    \label{itm:chapter5_a_shock_policies_cplt_pi_H}
        \(\pi_t^H\)：
        グロス・インフレ率\(\pi_t^H\)の図\ref{fig:chapter5_a_shock_graphs_pi_H}を見ると、CPLTの軌道は消費者物価インフレ目標（IT-CPI）と同様に、定常状態から極めて大きく不安定に乖離していることがわかる。
        総消費水準目標（CLT）が初期に大きなプラス方向へのスパイクを記録しているのと同様に、CPLTも供給ショックの直接的な影響を受けてショック直後にプラス方向（インフレ）へ鋭くスパイクしている。
        しかし、その後消費者物価を目標水準に引き戻そうとする強力な金融引き締めメカニズムが働くため、インフレ率は直ちに急反落し、極端なマイナス方向（デフレ）へと深く落ち込んでいる。この初期から中盤にかけてのインフレ率の激しい乱高下が次項で述べる価格分散の蓄積をもたらし、結果として厚生を悪化させる要因となっている。
    \item
    \label{itm:chapter5_a_shock_policies_cplt_Delta_t_minus_1_H}
        \(\Delta_{t-1}^H\)：
        価格分散\(\Delta_{t-1}^H\)の図\ref{fig:chapter5_a_shock_graphs_Delta_H}を見ると、CPLTの定常状態からの乖離は、生産者物価水準目標（PPLT）や生産者物価インフレ目標（IT-PPI）が全期間を通じてほぼゼロの水準を維持していることと比較して、明確な増大を見せていることがわかる。
        これは前述のインフレ率の初期における極端な落ち込みが価格分散として蓄積した結果であり、この生産効率の悪化がCPLTの厚生を押し下げる一因となっている。
        しかしながら、供給ショック（\(a\)ショック）においては価格分散の増大によるマイナス効果以上に、実体経済（消費や生産）の大暴落によるマイナス効果が極めて絶大である。そのため、CPLTの厚生が総消費水準目標（CLT）などに比べて低迷している主たる要因は、価格分散の蓄積よりもむしろ初期の深刻な不況を防げなかった点にあると言える。
\end{enumerate}